\section{Conclusion}
\label{sec:conclusion}

\pizf fails on rearranged scenes not because it cannot see where objects are, but because it does not act on it at the
moments that decide the task; on fixed layouts the training risk cannot tell the two apart.
Decision-time counterfactuals that move one variable, with labels valid after the move, identify use, and the change they
induce lives in the VLM's mid-to-late layers rather than in the action pathway.

\paragraph{Limitations.}
Our labels come from privileged scripted teachers in simulation; they cover reaching and placing, not articulated or
contact-rich steps, and their checks were refined by per-task failure analysis on the benchmark itself.
Results are for \pizf on one benchmark family with per-suite adaptation; gains on the multi-step suites are small, using the
scene costs robustness to appearance changes, and the early-approach blind spot shared by every model remains open.
